\chapter{Polynuclear magnetism (menu 39)}
\label{ch:menu39}
\menulabel{39}

\section{Purpose}

The POLY\_ANISO driver shipped inside the ORCA distribution (manual
section 7.18) combines the single-ion ab initio data of each magnetic
centre -- the \code{SINGLE\_ANISO} results of the \code{CASSCF/NEVPT2 + SOC}
runs, one data file per centre -- with the inter-site magnetic exchange and
the exact dipole-dipole coupling into the cluster magnetic properties.
This menu reads what the driver printed: the per-centre echo, the exchange
decomposition, the coupled states, $\chi T(T)$ and the Van Vleck
susceptibility tensors.  It never recomputes anything, and it never fits
the exchange constants.

\section{What you need}

A \code{poly\_aniso.output} produced by

\begin{quote}
\code{otool\_poly\_aniso < poly\_aniso.input > poly\_aniso.output}
\end{quote}

(the driver is called independently of ORCA), plus the per-centre data
files of the \code{SINGLE\_ANISO} runs (the
\code{<job>.CASSCF.\allowbreak anisofile} files) placed as
\code{aniso\_1.input}, \code{aniso\_2.input},
\ldots{} -- the names are mandatory (measured: a renamed file aborts the
run in \code{inquire\_key\_presence}).  The driver input starts with the
\code{\&POLY\_ANISO} head (equally mandatory -- without it the read stops
in \code{fetch\_init}) and carries the cluster definition: the \code{NNEQ}
block (centre types, equivalent counts, low-lying states per centre, with
the ab initio / isotropic-centre flags), the \code{COOR} coordinates and
the exchange constant list (\code{PAIR}/\code{LIN1}, \code{LIN3},
\code{LIN9}); the exchange constants are user input (measured elsewhere or
fitted), never computed by the program.

\section{Walkthrough}

The example reads the two-centre probe fixture (two Co(II) single-ion data
files at invented relative coordinates with an invented $J$ -- a format and
flow probe, see the boundaries).  The report echoes each centre (g values
2.0000 / 2.0000 / 1.9995 and the spin-orbit spectrum
0 / 0 / 29361.20 / 29361.20~cm$^{-1}$, matching the menu-36 analysis of the
same data), the exchange block ($J$ = 0.10000~cm$^{-1}$), the interaction
decomposition (LINES-1: isotropic 32.708\,\%, symmetric 93.839\,\%,
anti-symmetric 11.155\,\%), the coupled states and the $\chi T$ summary.

\lstinputlisting[style=fbkcode, firstline=57, lastline=84,
  title={The menu-39 example (the per-centre echo, the exchange
  decomposition, the coupled states and the $\chi T$ summary)}]
  {examples/expected/m39_polyaniso.txt}

The second leg of the same example writes the driver's input from the same
cluster description, with the checklist the workflow needs:

\lstinputlisting[style=fbkcode, firstline=146, lastline=164,
  title={The menu-39 input writer (mode 2): the questions and the written
  cluster summary}]
  {examples/expected/m39_polyaniso.txt}

\section{What you get}

The report \code{<output>.polyaniso.fbk.md}: the centre count; per centre
the data file, the coordinates, the spin-orbit state count and spectrum
and the g values; the exchange block (coupled states, pair list with the
$J$ values, which models were included); per pair and model the full
$3\times3$ interaction matrix with its isotropic / symmetric /
anti-symmetric decomposition weights; the coupled-state table (Lines /
dipole-dipole / total, absolute and relative); $\chi T(T)$ over the
requested temperature range and the Van Vleck susceptibility tensor
sequence (one $3\times3$ with main values and main axes per printed
temperature); the closing notes (skipped torque/magnetization by the
caller's choice, completion).

Mode 2 writes the driver's input instead: \file{<path>} (default
\file{./poly\_aniso.input}) carries the \code{NNEQ} block (the type count
with the ab initio flag, the equivalent centres and the spin-orbit
functions per type), the \code{PAIR} list (Lines-type
$\sum -J\,\mathbf{s}_i\cdot\mathbf{s}_j$), the optional \code{COOR} block
(symmetrised per-type coordinates; it switches the exact dipole-dipole
coupling on) and the optional \code{TINT} grid, closed by
\code{End of Input}; the report \file{<path>.fbk.md} repeats the cluster
summary, the data-file checklist (\file{aniso\_1.input}, \ldots), the run
command and the boundaries.

\section{Conventions, cross-checks and boundaries (measured)}

\begin{readbox}
\begin{itemize}
  \item \textbf{The exchange constants are the caller's input}: the driver
    (and this menu) never computes or fits them; a joint ab initio
    computation of exchange splittings (the LDF-CAHF / many-state
    PNO-CASPT2 route of Hallmen et al.\ 2019, a MOLPRO implementation) is
    outside the ORCA ecosystem -- registered as a documented termination.
  \item \textbf{The written input is accepted end to end}: the two-centre
    plan (2 types $\times$ 1 centre, spin-orbit basis $2+2$, one
    $J=0.1$~cm$^{-1}$ pair, the two coordinates, \code{TINT}
    $0$--$300$~K / $101$ points) drives \code{otool\_polyaniso} with the
    fixture's data files to rc\,=\,0 and reproduces the frozen probe output
    byte for byte (measured 2026-09-30).  The writer validates the
    structure (type count 1--6, at least one pair, pair indices,
    duplicate pairs, coordinate and grid shapes) but not the physics of
    the $J$ values or the coordinates; the \code{SYMM} and
    \code{LIN3}/\code{LIN9} input variants are registered.
  \item \textbf{The Lines model is exact only in three cases} (two
    isotropic spins; one Ising plus one isotropic spin; two Ising spins)
    and approximate otherwise; the dipole-dipole coupling is evaluated
    exactly from the ab initio moments and usually dominates in strongly
    anisotropic lanthanides (manual section 7.18).
  \item \textbf{Measured interface facts} (ORCA 6.1.1,
    \code{otool\_poly\_aniso} v1.0.0, fixtures \file{poly\_aniso/}): the
    \code{\&POLY\_ANISO} head and the \code{aniso\_N.input} names are
    mandatory; the driver output is recognised by the toolkit as
    ORCA-family output (its own banner); the per-centre echo cross-checks
    the menu-36 SINGLE\_ANISO analysis of the same data (the g and
    spectrum values agree).
  \item \textbf{Probe boundary}: the shipped fixture is a format-and-flow
    probe -- two independent single-ion runs at invented coordinates with
    an invented $J$ -- not a physical cluster; it validates the formats,
    the mandatory naming and the report parsing.
\end{itemize}
\end{readbox}
